\documentclass[10pt,a4paper]{amsart}
\usepackage[margin=19mm]{geometry}
\usepackage{amsmath,amssymb,mathtools}
\usepackage[scr=rsfs,cal=euler]{mathalfa}
\usepackage{xfrac}
\usepackage[unicode,hidelinks,bookmarks=false]{hyperref}
\newcommand{\ie}{i.e.\ }
\newcommand{\cf}{cf.\ }
\newcommand{\resp}{respectively\ }
\newcommand{\Subsection}{\subsection}
\providecommand{\nobreakdash}{\nobreak\hbox{-}}
\setlength{\emergencystretch}{2em}
\multlinegap0pt
\usepackage{amssymb}
\usepackage{tikz}
\newtheorem{lemma}{Lemma}
\newcommand{\Alt}{\mathrm{Alt}}
\newcommand{\C}{\mathbb{C}}
\newcommand{\DR}{\mathrm{DR}}
\newcommand{\Sym}{\mathrm{Sym}}
\newcommand{\mm}{\mathfrak{m}}
\newcommand{\sgn}{\mathrm{sgn}}
\newcommand{\xx}{\mathbf{x}}
\newcommand{\yy}{\mathbf{y}}
\title{Tautological classes for $(n,n+1)$ torus knots}
\author{Eugene Gorsky, Anton Mellit}
\date{Source: arXiv:2601.16877v1 (CC BY 4.0)}
\begin{document}
\maketitle

\noindent\emph{Source and licence.} Tautological classes for $(n,n+1)$ torus knots. Version arXiv:2601.16877v1 (CC BY 4.0), distributed by arXiv under the Creative Commons Attribution 4.0 International licence, http://creativecommons.org/licenses/by/4.0/, as recorded in the arXiv licence metadata for this exact version and shown on the abstract page https://arxiv.org/abs/2601.16877v1. Full licence notice: \texttt{licenses/arxiv-2601.16877-CC-BY-4.0.txt}.\par\smallskip

\noindent\textit{Editorial scope.} Counted unit: the article's lemma catalan (source0.tex lines 474--477). The article's complete original proof of it is delivered with it (source0.tex lines 479--496, one \texttt{proof} environment of 18 lines). No mathematics is written, repaired or re-derived: the statement and the proof are the article's own lines, in the article's own notation, and no retained line is substituted or re-flowed.  No further source-local context is needed: every cross-reference of every retained line resolves inside the two delivered spans.  Nothing was omitted from any retained span, so the package carries no omission record.  No retained line contains a citation, so the wrapper carries no bibliography and no bibliography entry.  No retained line contains a figure environment, an image-inclusion command or drawing code, so no image is delivered and the package declares no asset dependency.  Editorial formatting only, in addition: an AMS \texttt{amsart} preamble that restates the article's own declarations and macros used by the retained lines, namely the environments \texttt{lemma} on separate counters, together with the standard packages those lines need.  The retained environments print in this document as Lemma 1 in order of appearance, whereas the article prints the same items with the numbering of its own full document.  The complete native source of the article as distributed in the arXiv e-print for this exact version is bundled unchanged as \texttt{sources/193272/source0.tex}.
\begin{lemma}
\label{lem: catalan}
We have $\DR_n^{\sgn}\simeq J/\mm J$.
\end{lemma}

\begin{proof}
Let $L\subset \C[\xx,\yy]$ be the subspace spanned by the products $f(\xx,\yy)g(\xx,\yy)$ where $f$ and $g$ are homogeneous, $f$ has positive degree and symmetric, and $g$ is antisymmetric. Let 
$$
\Alt=\frac{1}{n!}\sum_{\sigma\in S_n}\sgn(\sigma)\sigma,\quad 
\Sym=\frac{1}{n!}\sum_{\sigma\in S_n}\sigma
$$
be the antisymmetrization and symmetrization operators in $\C[S_n]$ respectively.

Observe that any polynomial in $L$ is antisymmetric, and $L\subset \C[\xx,\yy]^{S_n}_{+}\cap \mm J$.  We claim that in fact $L$ is the sign isotypic component of both $\C[\xx,\yy]^{S_n}_{+}$ and $\mm J$. 
Indeed, any element in $\C[\xx,\yy]^{S_n}_{+}$ can be written as $\sum f_ih_i$ where $f_i$ are symmetric polynomials of positive degree. Then $\Alt(\sum f_ih_i)=\sum f_i\Alt(h_i)\in L$.
Similarly, any polynomial in $\mm J$ can be written as 
$\sum g_i h_i$ where $g_i$ are antisymmetric polynomials and $h_i$ are some homogeneous polynomials of positive degree. Then $\Alt(\sum g_ih_i)=\sum g_i\Sym(h_i)\in L$.

Since $J$ is generated by antisymmetric polynomials, we have
$$
J/\mm J=\C[\xx,\yy]^{\sgn}/(\mm J)^{\sgn}=\C[\xx,\yy]^{\sgn}/L=\DR_n^{\sgn}.
$$
\end{proof}

\end{document}
